\chapter{Conclusione}
Questo lavoro ha affrontato il problema del traffic shaping partendo dal suo funzionamento in ambienti Linux, 
per poi spiegarne i limiti in contesti containerizzati e, infine, introducendo due alternative interamente in user space. \\
\vspace*{0.01\textheight}
Inizialmente, sono stati spiegati i concetti fondamentali sul traffic shaping in ambiente, tramite il comando \texttt{tc},
spiegando le principali discipline di accodamento Hierarchical Token Bucket e Token Bucket Filter con scenari reali di utilizzo. \\ 
\vspace*{0.01\textheight}
Successivamente, è stato introdotto Kathará come ambiente di simulazione basato su container Docker, dimostrando, con esempi, il suo funzionamento.
In questa parte, sono stati analizzati i principali limiti della contaneirizzazione, spiegando come il comando \texttt{tc} necessiti di moduli kernel per poter funzionare correttamente. \\
\vspace*{0.01\textheight}
Dopo aver analizzato il problema, si è lavorato su un approccio diverso, spostando tutto la logica shaping del traffico dal kernel space in user space. \\
Il framework Netfilter, mettendo a disposizione le code NFQUEUE e il comando \texttt{iptables}, ha permesso di intercettare tutto il traffico dei processi dell'host, 
permettendo a un processo in stato utente di gestire la logica di shaping. In particolare, è stato scelto di usare il programma Damper, che tramite la libreria \texttt{libnetfilter\_queue}, 
implementa un algoritmo leaky bucket, per rallentare il traffico in uscita secondo diversi profili. \\
\vspace*{0.01\textheight}
Una volta dimostrata la fattibilità del nuovo approccio, si sono evidenziati i limiti strutturali sulla gestione di più flussi di traffico diversi, portando alla creazione di un suo fork: Donper.
In particolare, è stato modificato l'applicativo originale, supportando la gestione di più classi di traffico e di un algoritmo di condivisione di banda inutilizzata.
Inoltre, è stato sostituito l'utilizzo dell'algoritmo leaky bucket, in favore dell'algoritmo token bucket. \\
\vspace*{0.01\textheight}
Nel complesso, il lavoro dimostra che è possibile realizzare un traffic shaper, con funzioni simile a quelle offerte dal kernel Linux, interamente in userspace. \\ 
Restano, però, alcuni limiti: il programma Donper, supporta solamente un albero a due livelli, con un nodo padre unico. Un ulteriore aspetto da considerare riguarda la scalabilità dell'approccio al crescere della velocità delle interfacce di rete.
Le velocità di molte reti moderne, che possono raggiungere 2,5 o i 10 Gbit$/s$, potrebbero creare diversi problemi di bottleneck visto il numero di pacchetti al secondo da gestire. \\ 
\vspace*{0.01\textheight}
Un possibile sviluppo futuro potrebbe essere quello di valutare Donper in contesti più complessi, migliorando la gestione dei pacchetti, limitando il numero di context switch e sfruttando meglio l'utilizzo di thread.
